\section{Conclusion}
With the shown simulations of the both types of LMS algorithms, it easy to see how
strongly the output of the filter is affected by $\mu$ and $\rho$ as the step factors. 
Where a larger factor, in the case of the regular LMS, will result in a faster convergence, but
at the risk over a slight overshoot and a possible susceptability to sudden changes, whereas
small values will do the opposite, increase the time to convergence, in these tests, in the regular 
LMS, the increase was a factor of 10 each time, which showed that it had a larger impact, whereas 
having more filter tabs did not really impact the system, however as it is coefficients to calculates
it would inevitably have a price on computational speed.\\

However in the case of \textit{system identification}, it is possible to encounter a system 
which would require many more filter tabs, thus the need to have enought is a real one, and not 
having enough would lead to the inability to represent the actual system, thus residual errors would occur.\\

As for the variable step-size algorithm, it was possible to see how increasing the value of $\rho$ also
improved on the mean square error of signal, but at the cost of having the both slight overshoot
on the weights, which can seen on \ref{fig:VVSLMS_3_conv}